\chapter{The Residue Refuses To Disappear}

\section{The branch that cannot be silently discarded}

The residue capability consists of a step process and a replaceable representative question
\(\mathcal R(a,b)\).  For neutral accumulation forms
\(\operatorname{empty}(\rho)\) and \(\operatorname{indexed}(\rho,S,L)\), its default branch table is
\[
\begin{array}{c|c|c}
a&b&\mathcal R(a,b)\\
\hline
\operatorname{empty}(\rho_1)&\operatorname{empty}(\rho_2)&\rho_1=\rho_2\\
\operatorname{empty}(\rho_1)&\operatorname{indexed}(\rho_2,S_2,L_2)&\mathsf{yes}\\
\operatorname{indexed}(\rho_1,S_1,L_1)&\operatorname{empty}(\rho_2)&\mathsf{no}\\
\operatorname{indexed}(\rho_1,S_1,L_1)&\operatorname{indexed}(\rho_2,S_2,L_2)&
 [\rho_1=\rho_2]\land Q_{\le}(S_1,S_2)\land Q_{<}(L_1,L_2).
\end{array}
\]
The capability supplies no resolver for this question and proves no reflexivity, symmetry, transitivity, or equivalence
law.  Another declaration may replace \(\mathcal R\); every subsequent read must then name the replacement.

A history-bearing sample has the grammar
\[
\textit{sample}::=\operatorname{initial}(\rho,L)
\mid\operatorname{response}(\rho_s,S,\rho_r,R,H).
\]
The response stores a first receipt and stimulus, a second receipt and response, and prior history \(H\).  Its default
comparison question \(\mathcal C(x,y)\) has four branches:
\[
\begin{array}{c|c|c}
x&y&\mathcal C(x,y)\\
\hline
\operatorname{initial}(\rho_1,L_1)&\operatorname{initial}(\rho_2,L_2)&
 [\rho_1\sim\rho_2\land Q_{\le}(L_1,L_2)]
 \lor[\rho_1\nsim\rho_2\land Q_{\le}(L_2,L_1)]\\
\operatorname{response}(\rho_1,S_1,\rho'_1,R_1,H_1)&
\operatorname{response}(\rho_2,S_2,\rho'_2,R_2,H_2)&
 [\rho_1\sim\rho_2\land Q_{\le}(R_1,R_2)]
 \lor[\rho_1\nsim\rho_2\land Q_{\le}(S_1,S_2)]\\
\operatorname{initial}(\rho_1,L_1)&\operatorname{response}(\rho_2,S_2,\rho'_2,R_2,H_2)&
 [\rho_1\sim\rho_2\land Q_{\le}(L_1,S_2)]
 \lor[\rho_1\nsim\rho_2\land Q_{\le}(S_2,L_1)]\\
\operatorname{response}(\rho_1,S_1,\rho'_1,R_1,H_1)&\operatorname{initial}(\rho_2,L_2)&
 \mathsf{no}.
\end{array}
\]
Here \(\rho_1\sim\rho_2\) means derived receipt agreement.  In the response/response branch, agreement selects the
response comparison, while disagreement selects the stimulus comparison.  The second receipts
\(\rho'_1,\rho'_2\) and both history tails are ignored.  Again, the table is a question, not an established order.
The sample grammar supplies no resolver for \(\mathcal C\).

\paragraph{Worked example 1: receipt-selected comparison (computation).}
Take response samples \(X,Y\).  If their first receipts agree, the question reduces to
\(Q_{\le}(R_X,R_Y)\), regardless of their stimuli, second receipts, or histories.  If the first receipts disagree, it
reduces to \(Q_{\le}(S_X,S_Y)\), regardless of their responses.  Two samples can therefore receive the same comparison
answer while carrying different ignored histories.  Receipt relation selects an argument pair; it does not select
which pair has physical meaning.

Forward history construction prepends responses:
\[
\begin{aligned}
 H_0&=\operatorname{initial}(\rho_0,L_0),\\
 H_1&=\operatorname{response}(\rho_1,S_1,\rho'_1,R_1,H_0),\\
 H_2&=\operatorname{response}(\rho_2,S_2,\rho'_2,R_2,H_1),\\
 H_3&=\operatorname{response}(\rho_3,S_3,\rho'_3,R_3,H_2).
\end{aligned}
\]
Traversal reads newest to oldest: \(H_3,H_2,H_1,H_0\).

\paragraph{Worked example 2: budgeted history traversal (computation).}
With visit budget \(N=2\), traversal records the heads of \(H_3\) and \(H_2\), then stops with
\textsc{budget} and returns \(H_1\) as the residual prior tail.  Its certificate contains those two visited heads,
their receipts and payloads, the untouched tail \(H_1\), and the stop reason.  It makes no claim about the contents of
\(H_1\) or \(H_0\) beyond retaining their finite record.

Euclidean division is an external numerical model for explicit remainder accounting.  For \(n\ge0\), \(d>0\),
initialize \(q=0,r=n\).  While \(r\ge d\) and the budget remains, update
\[
                         q\leftarrow q+1,\qquad r\leftarrow r-d.
\]
The invariant is
\[
                         n=dq+r.
\]
A complete stop requires \(0\le r<d\); a budget stop preserves the invariant but does not certify a terminal Euclidean
remainder.

\paragraph{Worked example 3: iterative division (external numerical experiment).}
For \(n=23,d=5\), the trace is
\[
 (q,r):(0,23)\to(1,18)\to(2,13)\to(3,8)\to(4,3).
\]
Each row satisfies \(23=5q+r\).  The complete certificate gives \(q=4,r=3\), four transitions, the invariant checks,
\(d>0\), and \(0\le3<5\).  With budget two, the run stops at \((2,13)\).  It certifies
\(23=5\cdot2+13\), but not a completed division because \(13\ge5\).  This algorithm is an external arithmetic
experiment, not a property supplied by the residue capability.

Projection makes information loss measurable.  Define
\[
 \pi_{\rm now}(\operatorname{response}(\rho,S,\rho',R,H))=(\rho,S,\rho',R),
 \qquad
 \pi_{\rm full}(H)=H.
\]
The full projection is injective.  The current projection is many-to-one whenever histories differ behind an equal
head.

\paragraph{Worked example 4: four histories, two summaries (Shannon model).}
Let four histories \(H_1,H_2,H_3,H_4\) be equiprobable, with
\[
 \pi_{\rm now}(H_1)=\pi_{\rm now}(H_2)=u,\qquad
 \pi_{\rm now}(H_3)=\pi_{\rm now}(H_4)=v.
\]
Under this declared discrete distribution,
\[
 H(\text{history})=2\ \text{bits},\qquad
 H(\text{summary})=1\ \text{bit},\qquad
 H(\text{history}\mid\text{summary})=1\ \text{bit}
\]
by Shannon's definitions~\cite{shannon1948}.  The missing bit measures unresolved history within each two-member
fiber.  It is not heat, energy, or a physical erasure cost.

If \(H_A\ne H_B\) but \(\pi_{\rm now}(H_A)=\pi_{\rm now}(H_B)\), a separating question
\(q_{\rm full}:\mathsf{History}\to\mathsf{Answer}\) exists on the full finite history set, for example
\[
 q_{\rm full}(H)=[H=H_A].
\]
Then \(q_{\rm full}\circ\pi_{\rm full}\) distinguishes them.  For any
\(q_{\rm now}:\mathsf{Summary}\to\mathsf{Answer}\), no \(q_{\rm now}\) can separate equal summaries after
\(\pi_{\rm now}\), because both inputs have the same image.  Construction carries histories forward through either
projection; reading pulls the separately typed questions back:
\[
 \pi_{\rm full}^*q_{\rm full}=q_{\rm full}\circ\pi_{\rm full},
 \qquad
 \pi_{\rm now}^*q_{\rm now}=q_{\rm now}\circ\pi_{\rm now}.
\]

A checksum gives another lossy external projection.  For bytes \(b_1,\ldots,b_m\), let
\(c=(\sum_j b_j)\bmod7\).  Records \([1,2,3]\) and \([0,0,6]\) share checksum six while remaining distinct.  Equality
of checksums is not equality of records; the collision certificate identifies only their common projection.

\paragraph{Worked example 5: truncated acquisition (physics, then interpretation).}
Suppose an apparatus acquires five threshold-event records but a display retains three.  Computation can certify a
three-record prefix, two-record dropped suffix, and coverage \(3+2=5\).  Physics owns the sensor, threshold, acquisition
window, dead time, and missed-event model.  The displayed three is not a physical event count until that apparatus map
and its uncertainty are supplied.

The action ledger is explicit.  The residue capability, exact default representative and sample questions, history
updates, projections, budgets, and external arithmetic experiments are \emph{posited}.  The branch reductions,
ignored second receipts and histories, three-response trace, residual prior tail, division invariant, Shannon fiber
loss, checksum collision, and truncation coverage are \emph{demonstrated}.  Questions are \emph{read}
contravariantly through covariant projections and updates.  Order laws, numerical remainder meaning, physical residue,
and erasure energy remain \emph{reserved}.  A branch may be discarded intentionally, but the certificate must state
which branch, which fiber, and what information can no longer be recovered.

\section{Iterative error versus structural residue}

A state and a structural residue answer different questions.  The recurrence evolves states,
\[
                         x_{k+1}=T(x_k).
\]
A state error exists only relative to a declared reference \(x_*\), as \(e(x)=x-x_*\).
A structural residue is a retained branch, tail, receipt, or complement slice that a projection does not display.
The former may shrink, grow, alternate, or reach a fixed point.  The latter may have no numerical magnitude at all.
Calling both a remainder does not identify them.

For state \(x\), alleged target \(x_*\), update \(T\), finite approximation \(s_N\), exact target \(s_*\), and rounding
map \(Q\), the categories are
\[
\begin{array}{c|c|c}
\text{category}&\text{definition}&\text{required extra data}\\
\hline
\text{state error}&e(x)=x-x_*&x_*\\
\text{fixed-point residual}&r(x)=T(x)-x&T\\
\text{truncation error}&\tau_N=s_*-s_N&s_*\\
\text{rounding error}&\eta_Q(x)=Q(x)-x&Q\\
\text{structural residue}&\text{retained branch or projection fiber}&\text{record projection}
\end{array}
\]
A dropped suffix is not automatically numerical.  Without a valuation into an additive quantity, its numerical error
entry is \emph{absent}, not guessed.

Use the finite audit grammar
\[
\begin{array}{rcl}
\textit{state-run}&::=&
 \operatorname{run}(x_0,T,N,\textit{stop-question}),\\
\textit{loss-ledger}&::=&
 \operatorname{ledger}(\textit{stopping-loss},\textit{rounding-loss},
                       \textit{truncation-loss},\textit{retained-branch}),\\
\textit{stop}&::=&\textsc{budget}\mid\textsc{tolerance}\mid
 \textsc{fixed}\mid\textsc{resolver-domain}.
\end{array}
\]
For budget \(N\), compute \(x_{k+1}=T(x_k)\) for \(k=0,\ldots,N-1\).  The certificate contains \(x_0\), \(N\), all
transition pairs, the \(N+1\) visited states, declared projections and erasures, and stop reason.  State errors \(e\)
and fixed-point residuals \(r\) are computed ledger columns when their required data are supplied.  The certificate certifies no
unvisited update.

If an externally justified target \(x_*\) exists, a stopping report may include state error
\(d(x_N,x_*)\) under a declared distance.  For a finite record \(L\) sliced at \(K\), structural truncation loss is the dropped suffix
\(\operatorname{drop}_K(L)\), or a declared summary of it.  None of these is supplied merely by the recurrence.

Fixed-point and tolerance questions are reads on forward updates.  For
\[
 q_{\rm fix}(x)=[T(x)=x],\qquad
 q_\varepsilon(x)=[d(x,x_*)\le\varepsilon],
\]
the pulled-back questions are
\[
                         T^*q=q\circ T.
\]
The update is covariant; reading is contravariant.  A yes-answer to \(q_{\rm fix}\) certifies one fixed value.  A first
yes-answer to \(q_\varepsilon\) certifies one visited crossing.  Neither answer alone proves attraction, uniqueness, or
persistence.

\paragraph{Worked example 0: fixed residual, nonzero state error (external numerical experiment).}
Let \(T\) be the identity, \(x_*=0\), and \(x_0=7\).  Then
\[
                    r(x_0)=T(7)-7=0,\qquad e(x_0)=7-0=7.
\]
The update is at a fixed point, but that point is not the alleged target.  A zero fixed-point residual does not imply
zero state error without uniqueness or attraction premises.

\paragraph{Worked example 1: decreasing update error (external numerical experiment).}
Let \(x_0=8\), \(T(x)=x/2\), target \(x_*=0\), the ordinary norm, and \(N=4\).  Then
\[
 x_k=\frac{8}{2^k},\qquad e(x_k)=\frac{8}{2^k},\qquad
 r(x_k)=-\frac{4}{2^k},
\]
with trace \([8,4,2,1,\tfrac12]\).  If \(T\) is a contraction satisfying
\(\|T(x)-T(y)\|\le L\|x-y\|\), \(0\le L<1\), with fixed point \(x_*=T(x_*)\), then
\[
                  \|e(x)\|\le\frac{\|r(x)\|}{1-L}.
\]
For halving \(L=1/2\), so the bound is exact: \(\|e(x_k)\|=2\|r(x_k)\|\).
For tolerance \(\varepsilon=1\), the first accepted value is \(x_3=1\).  A first-hit policy stops there with
\textsc{tolerance}; a budget policy records one more update.  The trace demonstrates finite decrease.  Persistence
follows only after separately establishing \(0\le T(x)\le x\) for nonnegative \(x\).

\paragraph{Worked example 2: fixed magnitude without error reduction (external numerical experiment).}
Let \(T(e)=-e\), \(e_0=1\), and \(N=4\).  Then
\[
                         [1,-1,1,-1,1].
\]
The magnitude \(|e_k|=1\) is fixed while the signed error alternates.  A question about magnitude returns a constant
answer; a question about signed equality does not.  Selecting the former projection cannot justify a claim that the
full error reached a fixed point.

\paragraph{Worked example 3: nearest-hundredth ledger (external numerical experiment).}
Declare \(Q_2\) to round to the nearest hundredth, with exact halfway cases sent to the candidate whose final decimal
digit is even.  Then
\[
 Q_2(\tfrac13)=0.33=\frac{33}{100},\qquad
 \eta_{Q_2}(\tfrac13)=\frac{33}{100}-\frac13=-\frac1{300}.
\]
The certificate records exact input, rounded output, scale, tie rule, and signed error.  Changing the tie rule can
change halfway cases; the name rounding does not select a policy.

\paragraph{Worked example 4: truncation and retained complement (computation).}
For
\[
 L=[\mathtt a,\mathtt b,\mathtt c,\mathtt d,\mathtt e],\qquad K=3,
\]
stable slicing gives
\[
 \operatorname{keep}_3(L)=[\mathtt a,\mathtt b,\mathtt c],\qquad
 \operatorname{drop}_3(L)=[\mathtt d,\mathtt e].
\]
Coverage is \(3+2=5\).  The dropped slice is structural residue relative to the retained-prefix projection.  It is not
a scalar error until a separate valuation maps records to numbers.  Erasing the suffix destroys reconstruction; keeping
it in a complement ledger makes the loss explicit.

A history-bearing response provides the same distinction.  Let
\[
 H_A=\operatorname{response}(\rho,S,\rho',R,H_1),\qquad
 H_B=\operatorname{response}(\rho,S,\rho',R,H_2),\qquad H_1\ne H_2.
\]
The current projection maps both to \((\rho,S,\rho',R)\).  Their tails form a structural distinction even when the
current numerical error is equal.  A full-history question can separate them; no question on the common current summary
can.

\paragraph{Worked example 4b: projected fixed point, distinct histories (computation).}
Let a current-summary update be the identity.  The shared summary
\(z=\pi_{\rm now}(H_A)=\pi_{\rm now}(H_B)\) has projected residual
\[
                         T(z)-z=0.
\]
Nevertheless \(H_A\ne H_B\) because their tails differ.  The zero projected residual certifies stability of the
summary only; it does not eliminate the structural residue in the history fiber.

\paragraph{Worked example 4c: finite geometric truncation (external numerical experiment).}
For
\[
 s_N=\sum_{k=1}^{N}2^{-k}=1-2^{-N},
\]
declare exact target \(s_*=1\).  Then
\[
                         \tau_N=s_*-s_N=2^{-N}.
\]
The finite run owns the computed \(s_N\).  The symbolic error formula additionally depends on the declared target and
the separately established geometric identity; a finite prefix alone does not own an unrecorded infinite tail.

\paragraph{Worked example 5: finite information loss (Shannon model).}
Suppose four histories are equiprobable and a summary projection groups them into two equal fibers.  Then, under that
declared distribution,
\[
 H(\text{history})=2\text{ bits},\qquad
 H(\text{summary})=1\text{ bit},\qquad
 H(\text{history}\mid\text{summary})=1\text{ bit}
\]
by Shannon's definitions~\cite{shannon1948}.  The conditional bit quantifies unresolved history in this model.  It is
not a temperature, energy, or universal cost of erasure.

A physical error requires another provenance chain.  If a sensor report is \(y_k\) and a calibrated reference is
\(y_*\), physics may define \(e_k=y_k-y_*\) with units and uncertainty.  Computation can update, slice, and round the
finite recorded values.  It cannot infer that \(y_*\) is physically correct, that noise is independent, or that an
unobserved environment received discarded information.  Those are experimental premises.

The error ledger for a run is therefore
\[
 \mathcal E=(x_0,T,N,d,x_*,Q,K,\text{trace},\text{slices},
             e,r,\tau,\eta,\text{structural branches},\text{valuations},\text{stop}).
\]
Every unjustified entry is tagged \emph{absent} rather than guessed.  Every supplied target, norm, contraction
constant, rounding rule, and valuation is tagged \emph{external} or \emph{demonstrated}.  A budget stop reports the
last visited state and any error or residual columns justified by supplied reference data, together with the
unexecuted suffix of the request.  A tolerance stop reports its first accepted index.  A fixed stop reports the checked
pair \((x,T(x))\).  A truncation report records both slice sizes and the fate of the complement.

The action ledger closes the distinction.  Updates, metrics, targets, tolerances, contraction premises, rounding rules,
truncation budgets, valuations, and projections are \emph{posited}.  Finite traces, first hits, fixed pairs, residual
identities, contraction bounds under their premises, local rounding losses, geometric finite sums, coverage identities,
and projection fibers are \emph{demonstrated}.  Error and history questions are \emph{read}
contravariantly through covariant updates and projections.  Attraction, persistent enclosure, physical accuracy,
erasure energy, and numerical meaning for a structural branch remain \emph{reserved}.  Iterative error may be reduced;
structural residue must be accounted for even when it is not a number.

\section{Observation becomes history}

An observation capability stores a step process, records named before \(B\) and after \(A\), a relative question
\(Q_{\rm rel}\), and a sample update \(W\).  For accumulation forms
\(\operatorname{empty}(\rho)\) and \(\operatorname{indexed}(\rho,S,L)\), the relative question has the exact shape table
\[
\begin{array}{c|c|c}
B&A&Q_{\rm rel}(B,A)\\
\hline
\operatorname{empty}(\rho_1)&\operatorname{empty}(\rho_2)&R(\rho_1,\rho_2)\\
\operatorname{indexed}(\rho_1,S_1,L_1)&\operatorname{indexed}(\rho_2,S_2,L_2)&R(\rho_1,\rho_2)\\
\operatorname{empty}(\rho_1)&\operatorname{indexed}(\rho_2,S_2,L_2)&\mathsf{no}\\
\operatorname{indexed}(\rho_1,S_1,L_1)&\operatorname{empty}(\rho_2)&\mathsf{no}.
\end{array}
\]
Here \(R\) asks whether the first receipts agree in the declared sense.  The question ignores both chains and tails,
and the capability supplies no resolver.  Its name proves no temporal order, causality, or covariance.

Samples have the grammar
\[
 \textit{sample}::=\operatorname{initial}(\rho,L)
 \mid\operatorname{response}(\rho_s,S,\rho_r,R_s,\textit{sample}).
\]
With fixed receipt \(\rho_*\), the default observation update is
\[
\begin{aligned}
 W(\operatorname{initial}(\rho,L))
   &=\operatorname{response}(\rho,L,\rho,A,\operatorname{initial}(\rho,L)),\\
 W(\operatorname{response}(\rho_s,S,\rho_r,R_s,H))
   &=\operatorname{response}(\rho_s,S,\rho_*,A,
      \operatorname{response}(\rho_s,S,\rho_r,R_s,H)).
\end{aligned}
\]
The update retains the entire prior sample.  It preserves the first receipt and stimulus and replaces the current
response receipt and response.  It ignores the stored \(B\) and the stored step process.

For origin \(H_0\) and budget \(N\), compute \(H_{k+1}=W(H_k)\) for \(k=0,\ldots,N-1\).  The certificate contains all
transition pairs, \(N+1\) heads, retained receipts, declared erasures, and stop reason.  Tail traversal reads newest to
oldest and returns the unvisited prior tail on \textsc{budget}.

\paragraph{Worked example 1: accumulating observation history (computation).}
For \(H_0=\operatorname{initial}(\rho_0,L_0)\),
\[
\begin{aligned}
H_1&=\operatorname{response}(\rho_0,L_0,\rho_0,A,H_0),\\
H_2&=\operatorname{response}(\rho_0,L_0,\rho_*,A,H_1),\\
H_3&=\operatorname{response}(\rho_0,L_0,\rho_*,A,H_2).
\end{aligned}
\]
A traversal budget two records the heads of \(H_3,H_2\), returns \(H_1\), and certifies nothing about the unvisited
tail beyond retaining it.

A trial layer uses a different grammar:
\[
\begin{array}{rcl}
\textit{trial}&::=&\operatorname{hypothesis}(\rho,\textit{sample})\\
&&{}\mid\operatorname{trial-response}(\rho_1,\textit{sample},
                                      \rho_2,\textit{sample},\textit{trial}),\\
\textit{repeat-process}&::=&
 \operatorname{repeating}(\textit{observation-capability},
                          \textit{stored-stimulus},\textit{stored-expectation},P),\\
\textit{repeat-interface}&::=&
 \operatorname{repeatability}(\textit{repeat-process},\textit{typical-question}).
\end{array}
\]
Let \(S_*\) be the stored stimulus, \(E_*\) the stored expectation, and \(\rho_T,\rho_D\) fixed receipts.  The default
trial update is
\[
\begin{aligned}
 P(\operatorname{hypothesis}(\rho,S))
   &=\operatorname{trial-response}(\rho,S,\rho_T,S_*,E_*),\\
 P(\operatorname{trial-response}(\rho_1,S_1,\rho_2,S_2,T))
   &=\operatorname{trial-response}(\rho_2,S_*,\rho_D,S_2,E_*).
\end{aligned}
\]
The second branch ignores the first receipt \(\rho_1\), first sample \(S_1\), and prior tail \(T\).  Both branches reset
the prior tail to the stored expectation \(E_*\).  Thus \(W\) accumulates old samples, while \(P\) resets old trial
history.

\paragraph{Worked example 2: trial reset fiber (computation).}
Let
\[
 T_1=\operatorname{trial-response}(\rho_a,S_a,\rho_2,S_2,U_1),\qquad
 T_2=\operatorname{trial-response}(\rho_b,S_b,\rho_2,S_2,U_2),
\]
where the ignored first fields and tails may all differ.  Then
\[
            P(T_1)=P(T_2)=
            \operatorname{trial-response}(\rho_2,S_*,\rho_D,S_2,E_*).
\]
The update has a many-to-one fiber.  Equality after the reset does not reconstruct either prior trial.

The outer repeat interface supplies the default typical-response question
\[
                         \operatorname{typical}(s,t):=[P(s)=t].
\]
No resolver is supplied for it.  This unresolved question proves neither that two runs agree, that repeated applications return one
result, that a cadence is constant, nor that a laboratory experiment is repeatable.  Those require finite run
certificates and, for physics, an apparatus protocol.

Current and full-history projections remain distinct:
\[
 \pi_{\rm current}(\operatorname{response}(\rho_s,S,\rho_r,R_s,H))
       =(\rho_s,S,\rho_r,R_s),\qquad
 \pi_{\rm full}(H)=H.
\]
The first is many-to-one; the second is injective.  Construction moves histories forward through \(W\) or projections.
Questions are read contravariantly:
\[
 \pi_{\rm current}^*q_c=q_c\circ\pi_{\rm current},\qquad
 \pi_{\rm full}^*q_h=q_h\circ\pi_{\rm full}.
\]
No current-summary question separates equal summaries with different tails.

\paragraph{Worked example 3: current slice versus accumulated history (computation).}
Let two histories have equal current heads and unequal tails.  A current slice puts them in one fiber, while the
full-history question \(q_A(H)=[H=H_A]\) separates them.  With a traversal budget \(N\), coverage and exclusivity apply
only to the \(N\) visited rows; a stopped traversal makes no statement about its returned tail.

A sufficient statistic requires a declared probability family and target.  Let \(X,Y\) be independent uniform bits and
let the target be their parity
\[
                         \Theta=X\mathbin{\mathsf{xor}}Y.
\]
Use summary \(Z=\Theta\).  The four histories are equiprobable, so
\[
 H(X,Y)=2,\qquad H(Z)=1,\qquad H(X,Y\mid Z)=1
\]
in bits.  The summary discards one bit of history but is sufficient for the declared parity target because
\[
                         H(\Theta\mid Z)=0.
\]
By contrast, the first bit alone is not sufficient:
\[
                         H(\Theta\mid X)=1.
\]

\paragraph{Worked example 4: sufficiency is target-relative (Shannon model).}
The parity summary perfectly answers the parity question but cannot reconstruct which of the two histories in its fiber
occurred.  These entropy identities require the displayed joint distribution~\cite{shannon1948}.  Changing the target
or distribution requires a new sufficiency calculation.  The discarded bit is information in this model, not heat or
energy.

\paragraph{Worked example 5: apparatus history (physics, then interpretation).}
Suppose a detector creates time-stamped threshold records.  Physics owns the sensor, clock calibration, threshold,
dead time, missed-record rate, and uncertainty.  Computation can apply \(W\), \(P\), project the latest record, and
traverse finite history.  Neither prepend orientation nor a typical-response question establishes physical time,
constant cadence, causality, or experimental repeatability.

The observation ledger is
\[
 \mathcal O=(B,A,Q_{\rm rel},W,\text{repeat-process},P,\text{repeat-interface},
             S_*,E_*,N,\text{traces},
             \pi_{\rm current},\pi_{\rm full},\text{slices},
             \text{probability family},\text{target},\text{stop},\text{provenance}).
\]
Absent resolvers, distributions, apparatus maps, cadence models, and uncertainty are marked absent.

\paragraph{Descent experiment: a nonzero feature in a projection kernel.}
The Gibbs-preservation experiment gives structural residue an exact finite witness.
Its domain and codomain are both \(X=Y=\mathsf{Signal}=\mathsf{List}(\mathbb Z)\), and its transform groups adjacent samples by addition:
\[
 T_{\rm block}[x_0,x_1,x_2,x_3,\ldots]
   =[x_0+x_1,x_2+x_3,\ldots].
\]
The consumed base signal is \([2,4,6,8]\), while a feature of height \(h\) is \([h,-h]\), padded by pointwise signal addition.
For \(h=5\), the featured signal differs from the base signal but has the same block-sum shadow, because \((a+h)+(b-h)=a+b\).
Distinctness is an auxiliary theorem and a conjunct of the canonical top-level claim; the exported run decides only equality after projection.
Its run trace constructs the feature, adds it pointwise to the supplied base, applies recursive block summation to both signals, and checks list equality.
The commuting test is \(T_{\rm block}(S+F_h)=T_{\rm block}(S)\), and the residual test is \(S+F_h\ne S\) for \(h\ne0\).
No Fourier series or thermodynamic process is consumed; ``Gibbs preservation'' names only the witnessed structure
outside this shadow.

\paragraph{Descent experiment: message content forgets twist.}
The message experiment uses \(X=\mathbb Z^2\), \(Y=\mathbb Z\), and the receiver transform
\[
 T_{\rm msg}(a,b)=a+b.
\]
It separately defines the residual coordinate \(R(a,b)=a-b\) and the flip \(F(a,b)=(b,a)\).
The exact identities \(T_{\rm msg}\circ F=T_{\rm msg}\) and \(R\circ F=-R\) show that exchange leaves the received message fixed while reversing hidden twist.
The canonical input \((5,-5)\) and reference \((0,0)\) both project to zero, although the input residual is \(10\).
The run consumes only sent and reference updates and decides equality of their projected sums; changing the reference to \((3,3)\) returns false.
The trace consists of two additions and one integer equality, so it terminates immediately; its failure is unequal projected content, while equal content deliberately leaves the source update underdetermined.
The commuting test is flip invariance of \(T_{\rm msg}\), and the residual test is noninjectivity witnessed by the pure-twist and silent updates.
Electromagnetic curl, propagation, and a receiver circuit are absent; ``message'' and ``twist'' denote only the sum
and difference coordinates.

\paragraph{Descent experiment: two microhistories behind one modular record.}
The Butterfly-labelled experiment declares \(X=\mathsf{Microhistory}=\mathsf{List}(\mathbb N)\) and \(Y=\mathsf{CoarseRecord}=\mathsf{List}(\mathbb N)\).
For an apparatus cell size \(c>0\), its projection is
\[
 T_c[x_0,\ldots,x_n]=[x_0\bmod c,\ldots,x_n\bmod c].
\]
The canonical apparatus stores \(c=10\), capacity \(3\), and histories \([3,7,1]\) and \([13,17,21]\).
Those histories are unequal, yet both map to \([3,7,1]\), so the run reports them indistinguishable.
The exported run ignores capacity and compares only the two modular records.
A separate canonical theorem checks a separately defined three-history ledger against \texttt{defaultApparatus.capacity}; this capacity fact is a top-level claim conjunct, not a run-time horizon check.
The trace is one modular reduction per finite input cell followed by list equality; it stops when the lists end, and cell size zero would change the arithmetic meaning and therefore lies outside the stated \(c>0\) reading even if the language assigns a total remainder operation.
The representation test is pointwise agreement with map-by-remainder, while the residual consists of the discarded multiples of ten.
Butterfly dynamics, Lyapunov growth, and a sensor are absent; the chaos resonance is only underdetermination after
finite resolution.

\paragraph{Descent experiment: a balanced loop changes the path without changing balance.}
The first Gibbs experiment uses paths \(X=\mathsf{Ledger}(\mathsf{Step})\), with every step carrying an integer strain, and the balance map \(B:X\to\mathbb Z\) given by summation.
A catalyst with an integer pattern \([s_0,\ldots,s_n]\) generates the loop
\[
 L(s)=s\mathbin{++}(-\operatorname{reverse}s).
\]
The path transform \(T_s(P)=P\mathbin{++}L(s)\) consumes a bare path and the supplied catalyst.
The code proves \(B(L(s))=0\), \(B(T_s(P))=B(P)\), and \(\#T_s(P)=\#P+\#L(s)\).
For the canonical two-step bare path and four-step mirrored loop, the run returns count \(6\) while the loop balance evaluates to zero.
Its trace is finite reverse, sign negation, append, and fold; a nonmirrored catalyst is the explicit failure mode for
the zero-balance claim.
The commuting test is \(B\circ T_s=B\), whereas the residual test is the positive count difference showing that equal balance did not preserve the path.
Chemical Gibbs energy is absent; ``catalyst'' names only the supplied mirrored pattern.

The action ledger is exact.  The observation capability, trial grammar, inner repeat process, outer repeat interface,
default \(Q_{\rm rel},W,P\) and typical question, budgets,
projections, probability family, and target are \emph{posited}.  The ignored stored records and arguments, accumulating
observation trace, resetting trial fiber, finite slices, entropy identities, and target-relative parity sufficiency are
\emph{demonstrated}.  Current and history questions are \emph{read} contravariantly through covariant updates and
projections.  Temporal order, cadence, causal direction, whole-history recovery from a current summary, laboratory
repeatability, and sufficiency outside the declared model remain \emph{reserved}.  Observation becomes history under
\(W\); the later trial update \(P\) deliberately resets that history to its stored expectation.
